\section{问题提出的背景}

在智能制造、建筑仿真、航空航天等工业领域，高精度的计算机辅助设计（CAD）模型作为核心数据资产，通常具有极高的几何复杂度。近年来，随着工业设计与制造复杂性的急剧提升，模型的数量不断增长，对精度的需求也日益提高。这一发展趋势对传统的 CPU 驱动渲染管线提出了严峻挑战。

\subsection{背景介绍}

工业场景因其模型数量多且高度精细的特性，在实时渲染过程中面临诸多挑战：

在工业场景中，模型的多边形数量通常可达数亿级。传统的渲染方法要求将整个场景一次性加载至显存，即便在当前显存容量大幅提升的当下，仍可能引发显存溢出的问题。然而，实际上并不需要将整个场景同时存入显存，因为并非所有模型在任何时刻都是可见的。如果将这些不可见的模型加载至显存，不仅无法提升渲染效率，反而会导致大量显存资源的浪费。此外，对于距离相机较远的模型，也无需加载其最高精度的版本。由于这些模型在画面中所占比例较小，其细节损失对整体视觉效果影响有限，不会显著影响观感。

另一方面，在传统的 CPU 驱动渲染中， CPU 需承担场景遍历、剔除计算、渲染排序及绘制提交等任务。对于工业级场景而言，这种处理方式会导致 CPU 负载过高，进而严重制约 GPU 的高效运行\cite{pu2010}\cite{WangWei2011}。此外，工业场景通常包含数量庞大的对象，每个对象可能都会触发一次独立的绘制指令，从而显著降低渲染性能，导致帧率下降，难以满足实时交互的需求。

针对上述挑战，作为当前最成熟的商业引擎之一，虚幻引擎（Unreal Engine）提出了以 Nanite 为代表的 GPU 驱动渲染管线\cite{Nanite2022}。针对显存占用问题， Nanite 引入流式加载与虚拟化几何技术，实现按需动态加载资源，从而显著降低运行时的显存占用。针对渲染性能问题， Nanite 采用基于屏幕占用大小的动态 LOD（Level of Detail，细节层次）调整机制，有效优化渲染负载\cite{Overton2024}。此外，剔除计算、 LOD 管理等任务从 CPU 转移至 GPU 进行处理，不仅减少了 CPU 负担，还大幅提升了渲染性能，使得超高精度模型能够在保证视觉质量的同时，实现高效的实时渲染。

然而，由于历史遗留问题较多、系统架构复杂、受限于多平台兼容性，以及用户需求主要集中于游戏领域等因素，Nanite 在工业场景中的应用仍面临诸多挑战。具体而言，其预处理时间较长、帧率下降明显、内存占用严重等问题，限制了其在大规模工业数据可视化和实时渲染中的实际表现。

\subsection{本研究的意义和目的}

随着智能制造、建筑仿真、航空航天等工业领域的发展，高精度 CAD 模型的实时渲染与可视化需求日益增长。然而，现有的渲染方案主要针对游戏场景进行优化，在处理工业级大规模高精度数据时，往往面临显存占用过高、渲染性能下降、实时交互性不足等问题，难以满足工业应用的需求。因此，探索更适用于工业场景的轻量级渲染管线具有重要的现实意义。

本研究参考 Nanite 的核心思想，基于自研引擎设计并实现一套面向工业场景的 GPU驱动渲染管线，以提升大规模工业数据的实时渲染效率。研究重点包括优化流式加载策略以减少显存占用、改进 LOD 机制以适应工业模型的特点，并结合剔除计算以提高整体渲染性能。最终，研究成果将为工业级实时可视化、数字孪生、工程仿真等应用提供技术支撑，推动 GPU 驱动渲染在工业领域的进一步发展。